\chapter{Strong consistency and rates}
\label{chap:consistency}

This chapter proves the first main theorem: under the generic conditions and
Conditions~\textbf{A}--\textbf{B}, the allocation proportions converge almost
surely to the target and the estimators converge at the law-of-iterated-logarithm
rate. The proof is decomposed into small steps, each a candidate Lean lemma. It
corresponds to Theorem~3.1 and Appendix~A.2 of the paper.

Throughout the chapter we fix a RAR procedure satisfying the generic conditions
(Definition~\ref{def:generic_cond}) with some $\ell_{n,k}$, and we assume
Conditions~\textbf{A}--\textbf{B}. By Proposition~\ref{prop:aRTS_generic} this
covers every $\aRTS$ design.

\section{Convergence of the plug-in target}
\label{sec:cons_rho_limit}

\begin{lemma}[SLLN for the response martingale]\label{lem:respMart_div_count_slln}
\lean{AlphaRAR.respMart_div_pullCount_ae_tendsto_zero}\leanok
\uses{def:Q, def:counts, thm:mart_slln_bracket, lem:Q_quad_var}
For each arm $k$, on the event $\{N_{n,k}\to\infty\}$ the bracket-normalized response
martingale vanishes: $\dfrac{Q_{n,k}}{N_{n,k}}\to 0$ a.s.
\end{lemma}

\begin{proof}\leanok
\uses{thm:mart_slln_bracket, lem:Q_quad_var}
The response martingale $Q_{\cdot,k}$ has predictable quadratic variation
$\langle Q_k\rangle_n = V_k N_{n,k}$ (Lemma~\ref{lem:Q_quad_var}). If $V_k>0$, then on
$\{N_{n,k}\to\infty\}$ the bracket $\langle Q_k\rangle_n = V_k N_{n,k}\to\infty$, so the
bracket-normalized strong law (Theorem~\ref{thm:mart_slln_bracket}) gives
$Q_{n,k}/\langle Q_k\rangle_n\to0$, whence
$Q_{n,k}/N_{n,k} = V_k\,(Q_{n,k}/\langle Q_k\rangle_n)\to0$. If $V_k=0$, then
$\E[Q_{n,k}^2] = V_k\,\E[N_{n,k}] = 0$, so $Q_{n,k}=0$ a.s.\ for every $n$ and the ratio is
identically $0$.
\end{proof}

\begin{lemma}[Estimator limit on $\{N\to\infty\}$]\label{lem:theta_limit_infinite}
\lean{AlphaRAR.estimator_ae_tendsto_of_pullCount_atTop}\leanok
\uses{def:estimator, lem:estimator_bahadur, lem:respMart_div_count_slln}
For each arm $k$, on the event $\{N_{n,k}\to\infty\}$ the estimator converges a.s.\ to the arm
mean: $\hat\theta_{n,k}\to\theta_k$.
\end{lemma}

\begin{proof}\leanok
\uses{lem:estimator_bahadur, lem:respMart_div_count_slln}
By the exact error identity (Lemma~\ref{lem:estimator_bahadur}),
$\hat\theta_{n,k}-\theta_k = (Q_{n,k}+(\theta_{0,k}-\theta_k))/(N_{n,k}+1)$. On
$\{N_{n,k}\to\infty\}$ the constant offset is $O(N_{n,k}^{-1})\to0$, while
$\frac{Q_{n,k}}{N_{n,k}+1} = \frac{Q_{n,k}}{N_{n,k}}\cdot\frac{N_{n,k}}{N_{n,k}+1}\to0\cdot1=0$
by Lemma~\ref{lem:respMart_div_count_slln}. Hence $\hat\theta_{n,k}\to\theta_k$.
\end{proof}

\begin{lemma}[Estimator eventually constant on $\{\sup_n N<\infty\}$]\label{lem:theta_limit_finite}
\lean{AlphaRAR.exists_tendsto_estimator_of_not_pullCount_atTop}\leanok
\uses{def:estimator, def:counts}
For each arm $k$ and each realization, if $N_{n,k}\not\to\infty$ then $\hat\theta_{n,k}$ is
eventually constant, and in particular converges.
\end{lemma}

\begin{proof}\leanok
The count $N_{\cdot,k}$ is a nondecreasing $\N$-valued sequence, so $N_{n,k}\not\to\infty$
forces it to be bounded, hence eventually constant at its supremum $N_{N_0,k}$. For $n\ge N_0$
no further patient is assigned to arm $k$, so both the numerator $\sum_{j<n}X_{j,k}\xi_{j,k}$
and the denominator $N_{n,k}+1$ of $\hat\theta_{n,k}$ are frozen; hence $\hat\theta_{n,k}$ is
eventually constant.
\end{proof}

\begin{lemma}[The sequential estimator converges a.s.]\label{lem:theta_converges}
\lean{AlphaRAR.estimator_ae_tendsto}\leanok
\uses{lem:theta_limit_infinite, lem:theta_limit_finite}
For each arm $k$, the estimator $\hat\theta_{n,k}$ converges a.s.: to $\theta_k$ on
$\{N_{n,k}\to\infty\}$ and to an eventually-attained value on $\{\sup_n N_{n,k}<\infty\}$.
\end{lemma}

\begin{proof}\leanok
\uses{lem:theta_limit_infinite, lem:theta_limit_finite}
The alternative $N_{n,k}\to\infty$ versus $\sup_n N_{n,k}<\infty$ is exhaustive (as $N_{n,k}$
is monotone in $n$); branch~1 is Lemma~\ref{lem:theta_limit_infinite} and branch~2 is
Lemma~\ref{lem:theta_limit_finite}.
\end{proof}

\begin{lemma}[Dichotomy for the estimator limit]\label{lem:theta_limit_dichotomy}
\uses{def:estimator, def:condA, def:attainable, lem:theta_converges}
Under Condition~\textbf{A}, for each $k$ the estimator $\hat{\theta}_{n,k}$
converges a.s.\ to a limit in $I_k$: on $\{N_{n,k}\to\infty\}$ it converges to
$\theta_k$ by the strong law of large numbers; on $\{\sup_n N_{n,k} < \infty\}$
it is eventually constant and its limit lies in $V_k^{\mathrm{att}}$.
\end{lemma}

\begin{proof}
\uses{lem:theta_limit_infinite}
On $\{N_{n,k}\to\infty\}$, the estimator converges a.s.\ to $\theta_k$ by
Lemma~\ref{lem:theta_limit_infinite} (the martingale strong law applied to $Q_{\cdot,k}$). On
$\{\sup_n N_{n,k}<\infty\}$, only finitely many responses on arm $k$ are ever
observed, so $\hat{\theta}_{n,k}$ is eventually constant with a value in
$V_k^{\mathrm{att}}$.
\end{proof}

\begin{lemma}[The estimator vector converges a.s.]\label{lem:theta_converges_vec}
\lean{AlphaRAR.estimator_ae_tendsto_pi}\leanok
\uses{lem:theta_converges}
With finitely many arms, the per-arm a.s.\ limits of Lemma~\ref{lem:theta_converges} combine
into a single a.s.\ limit vector: almost surely there is $z\in\R^K$ with
$\estParams_n = (\hat\theta_{n,1},\dots,\hat\theta_{n,K}) \to z$.
\end{lemma}

\begin{proof}\leanok
\uses{lem:theta_converges}
Each component converges a.s.\ by Lemma~\ref{lem:theta_converges}; intersecting the finitely
many almost-sure events and choosing the (unique) limits componentwise gives a single a.s.\
event on which the whole vector converges.
\end{proof}

\begin{lemma}[Convergence of $\estTarget$ by continuity]\label{lem:rho_converges_cont}
\lean{AlphaRAR.rho_converges}\leanok
\uses{def:plugin_target, lem:theta_converges_vec}
For any continuous target map $\Target$, the plug-in target $\estTarget_n = \Target(\estParams_n)$
converges a.s.: almost surely there is $u\in\R^K$ with $\hat\rho_{n,k}\to u_k$ for every $k$,
where $u = \Target(z)$ and $z$ is the a.s.\ limit of $\estParams_n$.
\end{lemma}

\begin{proof}\leanok
\uses{lem:theta_converges_vec}
By Lemma~\ref{lem:theta_converges_vec}, $\estParams_n \to z$ a.s.\ (in the product topology of
$\R^K$). Continuity of $\Target$ transports this to
$\estTarget_n = \Target(\estParams_n) \to \Target(z) =: u$, and taking components gives
$\hat\rho_{n,k}\to u_k$.
\end{proof}

\begin{lemma}[Non-sparse limit of $\estTarget$]\label{lem:rho_converges}
\uses{def:condB, def:plugin_target, lem:rho_converges_cont}
Under Conditions~\textbf{A}--\textbf{B} there is a random $u \in (0,1)^K$ with
$\estTarget_n \to u$ a.s. Concretely $u = \Target(z)$ where $z\in I_1\times\dots\times I_K$
is the a.s.\ limit of $\estParams_n$.
\end{lemma}

\begin{proof}
\uses{lem:rho_converges_cont, lem:theta_limit_dichotomy}
By Lemma~\ref{lem:rho_converges_cont} (with $\Target$ continuous on $\hilb$ by Condition~\textbf{B})
and Lemma~\ref{lem:theta_limit_dichotomy}, $\estParams_n \to z$ a.s.\ with
$z \in I_1\times\dots\times I_K$. By Condition~\textbf{B}, $\Target$ is
continuous on the open set $\hilb \supseteq I_1\times\dots\times I_K$ and maps
$I_1\times\dots\times I_K$ into $(0,1)^K$; hence
$\estTarget_n = \Target(\estParams_n) \to \Target(z) =: u \in (0,1)^K$.
\end{proof}

\section{Convergence of the auxiliary process}
\label{sec:cons_U}

\begin{lemma}[LLN for the assignment martingale]\label{lem:M_lln}
\lean{AlphaRAR.assignMart_div_atTop_ae_tendsto_zero}\leanok
\uses{def:M}
$M_{n,k} = o(n)$ a.s.\ for each $k$.
\end{lemma}

\begin{proof}\leanok
\uses{lem:M_martingale, cor:mart_slln_bounded}
$(M_{n,k})$ is a martingale with bounded increments $|\Delta M_{m,k}|\le1$, so its
increments have summable normalized variances $\sum_m \frac{\Var(\Delta M_{m,k})}{m^2}<\infty$;
the strong law of large numbers for martingales gives $M_{n,k}/n \to 0$ a.s.
\end{proof}

\begin{lemma}[Limit of $U/n$]\label{lem:U_over_n}
\lean{AlphaRAR.auxU_div_tendsto}\leanok
\uses{def:U}
Under Conditions~\textbf{A}--\textbf{B}, for each $k$,
$\displaystyle \lim_{n\to\infty} \frac{U_{n,k}}{n} = -(1-\alpha)\,u_k$ a.s.,
where $u$ is the limit from Lemma~\ref{lem:rho_converges}. (The formalized statement
\texttt{auxU\_div\_tendsto} is the deterministic core: for real sequences with $\rho_n\to u$ and
$M_n/n\to0$, one has $U_n/n\to-(1-\alpha)u$; the a.s.\ statement is the pathwise application with the
a.s.\ limits from Lemmas~\ref{lem:M_lln} and~\ref{lem:rho_converges}.)
\end{lemma}

\begin{proof}\leanok
\uses{lem:M_lln, lem:rho_converges}
From Definition~\ref{def:U},
$\frac{U_{n,k}}{n} = \frac{\alpha}{n}\sum_{m=0}^{n-1}\hat{\rho}_{m,k}
 - \hat{\rho}_{n,k} + \frac{M_{n,k}}{n}$.
By Lemma~\ref{lem:M_lln} the last term is $o(1)$; by Lemma~\ref{lem:rho_converges}
$\hat{\rho}_{n,k}\to u_k$, and by Cesàro
$\frac1n\sum_{m<n}\hat{\rho}_{m,k}\to u_k$. Hence
$\frac{U_{n,k}}{n} \to \alpha u_k - u_k = -(1-\alpha)u_k$.
\end{proof}

\section{Matching of proportions and target}
\label{sec:cons_match}

\begin{lemma}[Positive part vanishes]\label{lem:pos_part_vanishes}
\lean{AlphaRAR.pos_part_vanishes, AlphaRAR.pos_part_vanishes_ae}\leanok
\uses{def:generic_cond}
Under the generic conditions and Conditions~\textbf{A}--\textbf{B}, for each $k$,
$\big(\frac{N_{n,k}}{n} - \hat{\rho}_{n,k}\big)^+ \to 0$ a.s. (The formalized statement
\texttt{pos\_part\_vanishes} is the pathwise core: it takes as hypotheses the plug-in convergence
$\rho_n\to u$ (\ref{lem:rho_converges}), the vanishing $M_n/n\to0$ (\ref{lem:M_lln}), and the two
generic conditions \eqref{eq:generic_ineq}--\eqref{eq:generic_small} in operational form.)
\end{lemma}

\begin{proof}\leanok
\uses{lem:convergence, lem:U_over_n}
Apply Lemma~\ref{lem:convergence} with $X_n := U_{n,k}$, $a_n := \ell_{n,k}$
(non-decreasing, $\le n$ by the generic conditions), and
$\varepsilon_n := \frac1n + \frac{N_{\ell_{n,k},k}-\ell_{n,k}\hat{\rho}_{\ell_{n,k},k}}{n}$,
which satisfies $\limsup\varepsilon_n \le 0$ by \eqref{eq:generic_small}. Since
$\frac{U_{n,k}}{n}\to -(1-\alpha)u_k$ (Lemma~\ref{lem:U_over_n}), the lemma yields
$\big(\varepsilon_n + \frac{U_{n,k}-U_{\ell_{n,k},k}}{n}\big)^+ \to 0$. Combined
with the generic inequality \eqref{eq:generic_ineq}, which bounds
$\frac{N_{n,k}}{n}-\hat{\rho}_{n,k}$ by that quantity, we get
$\big(\frac{N_{n,k}}{n}-\hat{\rho}_{n,k}\big)^+\to0$.
\end{proof}

\begin{lemma}[Negative part vanishes]\label{lem:neg_part_vanishes}
\lean{AlphaRAR.neg_part_vanishes, AlphaRAR.neg_part_vanishes_ae}\leanok
\uses{def:counts, def:plugin_target}
Under the same hypotheses, for each $k$,
$\big(\hat{\rho}_{n,k} - \frac{N_{n,k}}{n}\big)^+ \to 0$ a.s.
\end{lemma}

\begin{proof}\leanok
\uses{lem:pos_part_vanishes, lem:counts_sum}
Using $\sum_j N_{n,j}/n = 1$ (Lemma~\ref{lem:counts_sum}) and
$\sum_j \hat{\rho}_{n,j} = 1$,
$\big(\hat{\rho}_{n,k} - \frac{N_{n,k}}{n}\big)^+
 = \big(\sum_{j\ne k}(\frac{N_{n,j}}{n} - \hat{\rho}_{n,j})\big)^+
 \le \sum_{j\ne k}\big(\frac{N_{n,j}}{n} - \hat{\rho}_{n,j}\big)^+ \to 0$
by Lemma~\ref{lem:pos_part_vanishes}.
\end{proof}

\begin{lemma}[Proportions match the plug-in target]\label{lem:match}
\lean{AlphaRAR.match_proportion, AlphaRAR.match_proportion_ae}\leanok
\uses{def:counts, def:plugin_target}
Under the same hypotheses, for each $k$,
$\frac{N_{n,k}}{n} - \hat{\rho}_{n,k} \to 0$ and
$\frac{N_{n,k}}{n} \to u_k \in (0,1)$ a.s. (The formalized statement \texttt{match\_proportion}
delivers $\frac{N_{n,k}}{n}\to u_k$ from the two vanishing positive parts and $\hat{\rho}_{n,k}\to
u_k$.)
\end{lemma}

\begin{proof}\leanok
\uses{lem:pos_part_vanishes, lem:neg_part_vanishes, lem:rho_converges}
The absolute value $|\frac{N_{n,k}}{n}-\hat{\rho}_{n,k}|$ is the sum of the two
positive parts, both tending to $0$ by
Lemmas~\ref{lem:pos_part_vanishes}--\ref{lem:neg_part_vanishes}. Since
$\hat{\rho}_{n,k}\to u_k$ (Lemma~\ref{lem:rho_converges}), also
$\frac{N_{n,k}}{n}\to u_k$.
\end{proof}

\section{Identification of the limit and consistency}
\label{sec:cons_final}

\begin{lemma}[All arms sampled infinitely often]\label{lem:all_arms_infinite}
\lean{AlphaRAR.all_arms_infinite, AlphaRAR.all_arms_infinite_ae}\leanok
\uses{def:counts}
Under the same hypotheses, $N_{n,k}\to\infty$ a.s.\ for every $k$.
\end{lemma}

\begin{proof}\leanok
\uses{lem:match}
Since $\frac{N_{n,k}}{n}\to u_k > 0$ (Lemma~\ref{lem:match}), $N_{n,k}\sim u_k n
\to\infty$.
\end{proof}

\begin{lemma}[A.s. consistency of the proportions]\label{lem:consistency_matching}
\lean{AlphaRAR.consistency_ae}\leanok
\uses{lem:pos_part_vanishes, lem:neg_part_vanishes, lem:match, lem:rho_converges_cont}
Given the a.s.\ vanishing positive gaps for every arm (Lemma~\ref{lem:pos_part_vanishes}) and the
plug-in-target limits $\hat\rho_{n,k}\to u_k$ (Lemma~\ref{lem:rho_converges_cont}), together with
the simplex constraints ($\sum_k X_{j,k}=1$ and $\sum_k\hat\rho_{n,k}=1$), the proportions converge
a.s.\ to the target: $N_{n,k}/n \to u_k$ for all arms $k$ simultaneously.
\end{lemma}

\begin{proof}\leanok
\uses{lem:neg_part_vanishes, lem:match}
For each arm the negative gap vanishes by Lemma~\ref{lem:neg_part_vanishes}, and
Lemma~\ref{lem:match} then yields $N_{n,k}/n\to u_k$; intersecting the finitely many a.s.\ events
over the arms gives the simultaneous statement.
\end{proof}

\begin{lemma}[Estimator consistency]\label{lem:theta_consistent}
\uses{def:estimator, def:target, lem:consistency_matching}
Under the same hypotheses, $\estParams_n \to \Params$ a.s., and consequently
$\hat{\rho}_n \to v = \Target(\Params)$ and $\frac{N_{n,k}}{n}\to v_k$ a.s.
\end{lemma}

\begin{proof}
\uses{lem:all_arms_infinite, lem:theta_limit_dichotomy, lem:match}
By Lemma~\ref{lem:all_arms_infinite} every arm is sampled infinitely often, so
the first branch of Lemma~\ref{lem:theta_limit_dichotomy} applies to every $k$:
$\hat{\theta}_{n,k}\to\theta_k$. Continuity of $\Target$ gives
$\hat{\rho}_n\to\Target(\Params)=v$, and by Lemma~\ref{lem:match} the limit $u$ of
$N_n/n$ equals $v$.
\end{proof}

\section{Convergence rates}
\label{sec:cons_rates}

\begin{lemma}[LIL rate for the estimator]\label{lem:theta_LIL}
\lean{AlphaRAR.abs_estimator_sub_le_rate, AlphaRAR.abs_estimator_sub_le_rate_ae}\leanok
\uses{def:estimator}
Under the same hypotheses,
$\estParams_n - \Params = O\!\big(\sqrt{\tfrac{\log n}{n}}\big)$ a.s. (the $\log$, not $\log\log$,
form; see the LIL chapter's remark on the constant). The formalized statement
\texttt{abs\_estimator\_sub\_le\_rate} is the pathwise core: from the exact error identity
(Lemma~\ref{lem:estimator_bahadur}), a two-sided bound $|Q_{n,k}|\le C\sqrt{n\log n}$
(Lemma~\ref{lem:lil_truncation}) and $N_{n,k}/n\to v_k>0$ (Lemma~\ref{lem:match}), it gives
$|\hat\theta_{n,k}-\theta_k|\le C'\,\sqrt{n\log n}/n = C'\sqrt{\log n/n}$ for large $n$.
\end{lemma}

\begin{proof}\leanok
\uses{lem:all_arms_infinite, lem:Q_quad_var, lem:estimator_bahadur, lem:lil_truncation, lem:match}
By Lemma~\ref{lem:estimator_bahadur} the error is exactly
$\hat{\theta}_{n,k}-\theta_k = (Q_{n,k}+(\theta_{0,k}-\theta_k))/(N_{n,k}+1)$. The martingale
$Q_{\cdot,k}$ has quadratic variation $V_k N_{n,k}$ (Lemma~\ref{lem:Q_quad_var}), so the
finite-variance LIL (Lemma~\ref{lem:lil_truncation}, applied two-sided) gives
$|Q_{n,k}| = O(\sqrt{n\log n})$ a.s. Since $N_{n,k}\sim v_k n$ (Lemma~\ref{lem:match}), $N_{n,k}+1
\gtrsim (v_k/2)n$, and dividing bounds $|\hat{\theta}_{n,k}-\theta_k| \le
(2/v_k)(C+|\theta_{0,k}-\theta_k|)\sqrt{n\log n}/n = O(\sqrt{\log n / n})$.
\end{proof}

\begin{lemma}[Rate for the plug-in target]\label{lem:rho_rate}
\uses{def:plugin_target}
Under the same hypotheses,
$\hat{\rho}_n - v = O\!\big(\sqrt{\tfrac{\log\log n}{n}}\big)$ and hence
$n(\hat{\rho}_n - v) = O(\sqrt{n\log\log n})$ a.s.
\end{lemma}

\begin{proof}
\uses{lem:theta_LIL, def:condB, lem:taylor_rho}
A first-order Taylor expansion of $\Target$ around $\Params$ (twice
differentiable by Condition~\textbf{B}) gives
$\hat{\rho}_n - v = G(\estParams_n-\Params) + O(\|\estParams_n-\Params\|^2)$
with $G := \nabla\Target|_\Params$. Both terms are
$O(\sqrt{\log\log n / n})$ by Lemma~\ref{lem:theta_LIL}.
\end{proof}

\begin{theorem}[Strong consistency and rates]\label{thm:LLN}
\uses{def:generic_cond, def:condA, def:condB, def:counts, def:plugin_target}
Under the generic conditions and Conditions~\textbf{A}--\textbf{B}, for every
$k\in[K]$, as $n\to\infty$,
\[
  \frac{N_{n,k}}{n}\xrightarrow{a.s.} v_k, \qquad
  \estParams_n - \Params = O\!\Big(\sqrt{\tfrac{\log\log n}{n}}\Big)\ \text{a.s.},\qquad
  n(\hat{\rho}_n - v) = O(\sqrt{n\log\log n})\ \text{a.s.}
\]
\end{theorem}

\begin{proof}
\uses{lem:theta_consistent, lem:theta_LIL, lem:rho_rate}
Combine Lemma~\ref{lem:theta_consistent} (first statement), Lemma~\ref{lem:theta_LIL}
(second) and Lemma~\ref{lem:rho_rate} (third).
\end{proof}
